\documentclass[11pt]{article}

\usepackage{amsmath,amssymb,amsthm}
\usepackage[margin=1in]{geometry}
\usepackage{xcolor}
%\usepackage{url}
%\usepackage{float}
%\usepackage{array}
\usepackage{comment}
\usepackage{booktabs}
\usepackage{graphicx}
\usepackage{tcolorbox}
%\usepackage{tikz}
%\usetikzlibrary{decorations.pathreplacing,arrows.meta,calc,patterns}
\usepackage{enumitem}
%\usepackage{csvsimple}
\usepackage{hyperref}
\usepackage[capitalize,sort]{cleveref}
%\usepackage{algorithm,algpseudocode}

% colors
\colorlet{textColor}{black}
\colorlet{bgColor}{white}
\colorlet{textRed}{red!50!black}
\colorlet{textGreen}{green!50!black}
\colorlet{textBlue}{blue!50!black}
\colorlet{textHeavy}{black}
\colorlet{dimColor}{textColor!50!bgColor}

\hypersetup{
colorlinks,linkcolor=textRed,citecolor=textRed,urlcolor=textBlue,
hypertexnames=false,bookmarksnumbered=true,final
}

% macros
\newcommand*{\Th}{^{\textrm{th}}}
\newcommand*{\abs}[1]{\lvert #1 \rvert}
\newcommand*{\defeq}{:=}
\DeclareMathOperator*{\E}{\mathbb{E}}
\DeclareMathOperator*{\Var}{Var}
\DeclareMathOperator*{\argmin}{argmin}
\DeclareMathOperator*{\argmax}{argmax}
\renewcommand*{\Pr}{\operatorname{\mathrm{P}}}

% theorems, lemmas, definitions
\theoremstyle{plain}
\newtheorem{theorem}{Theorem}
\newtheorem{lemma}[theorem]{Lemma}
\theoremstyle{definition}
\newtheorem{definition}{Definition}
\newtheorem{example}{Example}
\newtheorem{remark}{Remark}

% wrap URLs
\Urlmuskip=0mu plus 0.1mu
\makeatletter
\g@addto@macro{\UrlBreaks}{%
\do\/%
\do\a\do\b\do\c\do\d\do\e\do\f\do\g\do\h\do\i\do\j\do\k\do\l\do\m%
\do\n\do\o\do\p\do\q\do\r\do\s\do\t\do\u\do\v\do\w\do\x\do\y\do\z%
\do\A\do\B\do\C\do\D\do\E\do\F\do\G\do\H\do\I\do\J\do\K\do\L\do\M%
\do\N\do\O\do\P\do\Q\do\R\do\S\do\T\do\U\do\V\do\W\do\X\do\Y\do\Z%
\do\0\do\1\do\2\do\3\do\4\do\5\do\6\do\7\do\8\do\9%
}
\makeatother

% tighter lists
\newenvironment*{tightemize}{\begin{itemize}[noitemsep]}{\end{itemize}}%
\newenvironment*{tightenum}{\begin{enumerate}[noitemsep]}{\end{enumerate}}%

% customizations
\allowdisplaybreaks
%\defaultaddspace=0.0em  % spacing in tables

\newcommand*{\bibPrep}{
\phantomsection
\addcontentsline{toc}{section}{References}
\bibliographystyle{plainurl}
}


\newcommand{\lori}{\ell_{\mathrm{ori}}}
\newcommand{\ldes}{\ell_{\mathrm{des}}}

\tcbuselibrary{theorems, skins, breakable}
\newtheorem{exercise}{Exercise}
\crefname{exercise}{Exercise}{Exercises}

\tcbset{
exeStyle/.style={colback=textColor!3!bgColor, frame empty, colupper=textColor},
}
% \newtcbtheorem[number within=section]{exercise}{Exercise}{mytheoremstyle}{exe}

\title{IE 300 Case Study 1: Fastest Airlines}
\author{\empty}
\date{\empty}

\begin{document}

\maketitle
%\setlength{\parindent}{0pt}
\setlength{\parskip}{0.4em}

%\section{Learning Objectives}

This study aims to help you develop basic Python skills for data analysis.
You will learn how to assess airlines based on historical flight data.
%
By the end of this case study, you will be able to:
\begin{tightemize}
\item Read data from csv files.
\item Work with different data structures in Python.
\item Clean data and perform calculations.
\item Create basic visualizations using \texttt{matplotlib}.
\end{tightemize}

\section{How do we determine which airline is the ``fastest''?}

One can define the \emph{duration} of a flight to be the difference between
the scheduled arrival and departure times (after adjusting for timezone changes, if any),
and compare flights across the same route based on these durations.
But these durations are not always accurate. Airlines don't want to be labeled \emph{late},
so they often pad their schedules, and so, report longer flight times.

A different approach is to record \emph{actual} arrival and departure times, instead of scheduled times.
But actual times can vary, because of weather, airport congestion, etc.
So, given historic flight data, we would like to find how long a flight takes \emph{on average}.

The dataset \texttt{flight-time.csv}, from the Bureau of Transportation Statistics,
contains 743 observations of flights from ORD to LAX throughout November 2015.
There are 10 columns in \texttt{flight-time.csv}, described in \cref{table:ft}.

\begin{table}[htb]
\centering
\caption{Columns in \texttt{flight-time.csv}}
\label{table:ft}
\begin{tabular}{cll}
\toprule
\# & Column Name & Column Description
\\ \midrule
1 & Flight Date & Date of the flight (\texttt{mm/dd/yyyy})
\\ 2 & Carrier & IATA Carrier Identification
\\ 3 & Flight Number & Number assigned to the flight
\\ 4 & Origin & Origin airport
\\ 5 & Destination & Destination airport
\\ 6 & Departure Time & Time of departure (\texttt{hhmm})
\\ 7 & Departure Delay & Difference between scheduled and actual departure time (minutes)
\\ 8 & Arrival Time & Time of arrival (\texttt{hhmm})
\\ 9 & Arrival Delay & Difference between scheduled and actual arrival time (minutes)
\\ 10 & Flight Time & Difference between actual departure and arrival time (minutes)
\\ \bottomrule
\end{tabular}
\end{table}

There are four types of time we are interested in:
\begin{enumerate}
\item \emph{Average Flight Time}:
    The average time an airline takes to complete the route.
\item \emph{Target Flight Time} (TFT):
    An estimate of how long a flight should take based on distance and direction of travel.
    For this case study, we will use a simple formula that does not
    include complex variables like wind speed, jetstreams, etc.
    We will calculate it based on the spherical distance (great circle distance).
    Let $\lori$ and $\ldes$ be the longitudes of the origin and destination, respectively,
    and $d$ be the spherical distance between two points (in miles).
    Assuming a constant velocity and an average time of 20 minutes to runway time,
    define the target flight time (in minutes) as
    \[ \mathrm{TFT} = 0.117 \times d + 0.517 \times (\lori - \ldes) + 20. \]
\item \emph{Typical Time}:
    The target flight time, plus the average delay associated with the origin airport,
    plus the average delay associated with the destination airport.
\item \emph{Time Added}:
    The measure of an airline's performance is the average flight time minus the typical time.
    This difference is called \emph{time added}.
    The lower the time added, the faster the airline is on that particular route.
\end{enumerate}

\begin{tcolorbox}[exeStyle]
\begin{exercise}
\label{exe:1}
Write Python code to do the following:
\begin{enumerate}[label=(\alph*)]
\item Read the dataset (\texttt{flight-time.csv}), ignoring any rows without a recorded departure or arrival time.
    Also, some times were recorded incorrectly, so ignore rows where the flight time is less than 230 minutes.
    Calculate and print the number of remaining observations in the dataset.
\item Calculate and print the target flight time with $d = 1741.16$ miles,
    $\lori = -87.90^{\circ}$ and $\ldes = -118.41^{\circ}$.
    (These values are for the route from ORD to LAX.)
\item Calculate and print the typical time of this route.
\item Calculate and print the time added for each airline.
    Which airline has the lowest time added?
\end{enumerate}
\end{exercise}
\end{tcolorbox}

\begin{tcolorbox}[exeStyle]
\begin{exercise}
\label{exe:2}
Using \texttt{matplotlib}, create a bar graph for the time added of each airline.
Be sure to label your axes and the plot.
\end{exercise}
\end{tcolorbox}

\nocite{nate2015how}
\bibPrep{}  % \bibPrep is defined in header.tex
\bibliography{src/bibdb}

\newpage
\appendix
\section{Grading Rubric (40 points)}

\subsection{Exercise \ref{exe:1} (30 points)}

\begin{enumerate}[label=(\alph*)]
\item Read and clean the data (12 points):
    \begin{tightemize}
    \item 3 points to correctly read \texttt{flight-time.csv}.
    \item 3 points to remove rows with missing data.
    \item 3 points to remove rows with incorrect flight time.
    \item 3 points to print the number of observations.
    \end{tightemize}
\item Target flight time (4 points).
\item Typical time (5 points).
\item Time added and best airlines (9 points):
    \begin{tightemize}
    \item 6 points to compute the time added for each airline.
    \item 3 points to print the best airlines.
    \end{tightemize}
\end{enumerate}

\subsection{Exercise \ref{exe:2} (10 points)}

\begin{tightemize}
\item 6 points for the bar graph.
\item 2 points to label the axes.
\item 2 points to label the plot.
\end{tightemize}

\end{document}
